\section{The Hilbert Door}

Energy squared is enough to detect activity, but it is not yet the whole
language of scale. A square is excellent for positivity. It refuses negative
activity and, after anchoring, it vanishes only when the activity itself
vanishes. But a reader also needs a magnitude that can be compared with
other magnitudes, transported through refinement, and used to say how far a
finite state has moved. The square root supplies that magnitude.

The door has the simple shape
\[
  \|x\| = \sqrt{E(x)^2}.
\]
Again the formula is not the authority. The authority is the grammar that
has made the formula honest. The anchored chain has already forced
\(E(x)^2=0\) to identify zero activity. The square root now turns the
positive quadratic reading into a scale reading. It carries the same zero
detection, but in the units of magnitude rather than squared activity.

This matters because refinement is not always read directly in squares. A
finite chain may be compared to another finite chain. A partition may be
split. A supported correction may be added. The grammar needs to say when
these changes are small, when they are shrinking, and when two refinement
histories have become indistinguishable by every permitted activity test.
The norm gives that language. It reads activity as distance inside the
finite grammar.

The Hilbert door opens when finite activity becomes complete enough to house
limits. But the door opens only after the zero test has been earned. If
\(\|x\|=0\) could hold for a nonzero admissible activity, completion would
identify distinct readings by mistake. If a refinement sequence could have
shrinking pairwise differences while hiding an unanchored drift, the
completed point would inherit the ambiguity. The anchor and the discrete
Poincare reading therefore stand at the threshold. They make the norm a real
scale rather than a quotient still leaking null motion.

Completion should not be confused with possession of every point in advance.
The grammar does not begin with a completed Hilbert space and then decide to
measure inside it. It begins with finite chains whose activity can be read.
It permits refinement when the activity of the difference between stages is
bounded and forced downward. It identifies refinement histories when their
difference has norm zero. The completed place is the home earned by those
finite histories, not a hidden room that was available before the work began.

Sampling gives the right warning. A sequence of recorded samples can support
exact recovery only when the ledger has committed enough distinctions to
make recovery lawful. A smooth-looking reconstruction is not justified by
appearance. It is justified when the recorded refinement carries enough
activity information to forbid two different completions from reading the
same finite data. The norm plays the same role here. It forbids completion
from pretending that unresolved activity has disappeared.

This is why the square-root door is also a scale-field door. The real scale
field is the range in which activity magnitudes can be compared. A node, a
chain, a label, or a later sensor reading may carry many kinds of
presentation, but their activity must be measured in one scale if they are
to be compared as refinements of the same invariant. The norm provides that
scale. It lets the grammar read smallness, convergence, and zero activity
without abandoning the finite origin of the record.

The square root is not a decorative operation at the end of a calculation.
It changes what the invariant can do. Energy squared says how much
quadratic activity is present. Its square root says how large the activity
is as a readable magnitude. The first protects positivity. The second
permits metric comparison. Together they let the concrete operator carry
the invariant into a completed scale without losing the boundary discipline
that made positivity possible.

The door also clarifies the relation between cost and burden from the first
section. Elapsed cost may be counted by a clock, but the size of a correction
must be read by a scale. A resolution that spends many small refinements can
converge; a resolution that spends fewer but larger unresolved activities
may not. The grammar needs both counts and magnitudes. The clock reads
elapsed turns; the norm reads the activity each turn still carries.

Once the scale field is available, the invariant can be transported beyond
the bare chain. It can attach to labels, support phases, and sort kinds of
readings without being replaced by them. The next problem is therefore not
how to invent more invariants. It is how one invariant can bind the labels
by which a universe becomes readable.
